\documentclass[12pt,a4paper]{article}
%per compilare il template usare questo sotto e commentare l'altro:
%\usepackage{../stile}
\usepackage{../../../template/stile}


%Titolo documento
\newcommand{\titoloDocumento}{Norme di Progetto}

%Prima data di creazione del documento
\newcommand{\dataCreazione}{30 novembre 2015}

%Inserite la versione attuale del documento
\newcommand{\versione}{1.0.0}

%Stato in cui si trova il documento: Formale solo all'atto di consegna
\newcommand{\stato}{Informale}

%Uso del documento
\newcommand{\uso}{Interno}

\rhead{\titoloDocumento}
\lfoot{Versione: \versione}
\title{\titoloDocumento}

\begin{document}
\begin{titlepage}
\begin{center}
\today \\
\vspace{1cm}
\begin{Huge}
\textbf{\nomeTeam}
\end{Huge}
\vspace{1cm}
\includegraphics[scale=0.5]{\logoGrande}
\vspace{1cm}

\HRule \\[0.4cm]
\begin{Huge}
{\huge \bfseries \titoloDocumento}\\[0.4cm]
\end{Huge}
\HRule \\[1cm]
\vfill

\begin{table}[h]
\begin{center}
\begin{tabular}{r | l}
\multicolumn{2}{c}{\textbf{Informazioni sul documento}}\\
\midrule
\textbf{Nome Documento}	&	\titoloDocumento	\\
\textbf{Versione}	&	\versione	\\
\textbf{Stato}	&	\emph{\stato}	\\
\textbf{Uso}	&	\emph{\uso}	\\
\textbf{Data Creazione}	&	\dataCreazione	\\
\textbf{Data Ultima Modifica}	&	\today	\\
\textbf{Redazione}	&	\NDC	\\
\ &	\AVE	\\
\ &	\AVI	\\
\textbf{Verifica}	&	\IB	\\
\ & \TP \\
\textbf{Approvazione}	&	\AB	\\
\textbf{Lista Distribuzione}	&	\nomeTeam	\\

\end{tabular}
\end{center}
\end{table}

\end{center}
\end{titlepage}
\newpage

\Large{\textbf{Registro delle modifiche}}\\
\normalsize

\begin{table}[h]
\begin{center}

\begin{tabular}{p{0.12\textwidth} p{0.25\textwidth} p{0.18\textwidth} p{0.4\textwidth}}
\toprule
\textbf{Versione}	&	\textbf{Autore}	&	\textbf{Data}	&	\textbf{Descrizione}\\
\midrule
\midrule
1.0.0 & \NDC & 2015-11-30 & Template di base \\
\bottomrule
\end{tabular}
\caption{Versionamento del documento}
\label{tabVers1}
\end{center}
\end{table}
\newpage

\tableofcontents
\newpage

\listoftables
\listoffigures
\newpage

% NOTA non so dove mettere l'uso di git

\section{Introduzione}
Questo documento definisce le norme che i membri del gruppo \nomeTeam{} adotteranno nello svolgimento del progetto \CapitolatoL. Tutti i membri sono tenuti a leggere il documento attentamente e a seguirene le norme.

\subsection{Scopo}
Il documento si propone come proposito di garantire l’uniformità di tutto il materiale prodotto, migliorarene l’efficienza e ridurrene gli errori.

\subsection{Descrizione}
Verranno definite norme riguardanti:

\begin{itemize}
  \item le interazioni tra membri del gruppo;
  \item la stesura documenti, convenzioni stilistiche e tipografiche;
  \item le modalità di lavoro durante tutte le varie fasi del progetto;
  \item l'ambiente di lavoro.
\end{itemize}

\subsection{Glossario}
Al fine di evitare ogni ambiguità di linguaggio e massimizzare la comprensione dei
documenti, i termini tecnici, di dominio, gli acronimi e le parole che necessitano chiarimenti sono riportate nel documento \emph{Glossario\_v1.2.0}. Ogni occorrenza di vocaboli presenti nel \emph{Glossario} è marcata da una G maiuscola in pedice (e.g. pedice\textsubscript{G}).

\newpage

\section{Processi di Sviluppo} %ISO 12207 5

\subsection{Scopo}

\subsection{Descrizione}

\subsection{Analisi dei Requisiti} %ISO 12207 5.3.2

\newpage

\section{Processi di Supporto} %ISO 12207 6

\subsection{Scopo}
Definisce norme per lo sviluppo e il mantenimento della documentazione prodotta durante il ciclo di vita del software. Inoltre definisce metodi per il controllo della qualità, di verifica e validazione di tali documenti.


\subsection{Descrizione}
Verranno definite norme riguardanti:

\begin{itemize}
  \item l'utilizzo e l'accesso ai documenti;
  \item i templates da utilizzare, l'uniformità di linguaggio, le convenzioni stilistiche e tipografiche;
  \item i metodi di verifica e approvazione dei documenti;
\end{itemize}

\subsection{Documentazione} %ISO 12207 6.1

Per agevolare la redazione della documentazione è stato creato un template \LaTeX{} contenente tutte le impostazioni stilistiche e grafiche citate in questo documento. Tale modello si può trovare nella repository\textsubscript{G} in \texttt{doc/template}.

Nei documenti alcuni termini devono essere interpretati in modo analogo ai termini inglesi descritti in RFC 2119\footnote{http://www.ietf.org/rfc/rfc2119.txt}:

\begin{itemize}
  \item i termini \emph{deve}, \emph{è richiesto} e sinonimi stretti sono da intendersi con lo stesso significato di \emph{MUST};
  \item i termini \emph{non deve}, \emph{è richiesto che non} e sinonimi stretti sono da intendersi con lo stesso significato di \emph{MUST NOT};
  \item i  termini \emph{dovrebbe}, \emph{si raccomanda}, \emph{è preferibile} e sinonimi stretti sono da intendersi con lo stesso significato di \emph{SHOULD}
  \item i termini \emph{non dovrebbe}, \emph{si raccomanda di non}, \emph{è preferibile che non} e sinonimi stretti sono da intendersi con lo stesso significato di \emph{SHOULD NOT};
  \item i termini \emph{può}, \emph{opzionalmente} e sinonimi stretti sono da intendersi con lo stesso significato di \emph{MAY}.
\end{itemize}

\subsection{Configurazione} %ISO 12207 6.2

\subsection{Verifica} %ISO 12207 6.4

\subsection{Validazione} %ISO 12207 6.5

\subsection{Risoluzione dei Problemi} %ISO 12207 6.8

\newpage

\section{Processi di Organizzazione} %ISO 12207 7.1

\subsection{Scopo} %ISO 12207 7.1.1

\subsection{Descrizione} %ISO 12207 7.1.1

\subsection{Pianificazione} %ISO 12207 7.1.2
%task ruoli e responsabilità

\subsection{Esecuzione e controllo} %ISO 12207 7.1.3
%come viene monitorata l'esecuzione della pianificazione: DA DEFINIRE (per esempio chiedendo ogni tanto come sta andando il lavoro)

\subsection{Infrastruttura} %ISO 12207 7.2
%asana per i task, github per le issues e bug e risoluzione problemi, projectlibre per la pianificazione con gantt e costi

\newpage


\end{document}
